%% sections/threats.tex -- Phase 7.1. Register: CLAUDE.md S9.2. REWRITTEN per PHASES 7.1: the old
%% "two-teacher consensus, human-validated" wording is the S2.3 over-claim and must not reappear.
%% Groups: gold/data (reliability, criterion limit, inflation-as-documented-not-supported),
%% measurement/reproducibility, scope of the claims, data/code availability.
%% COI statement REMOVED from the paper + dropped as a CLAUDE.md non-negotiable (user decision 2026-06-17):
%% collecting one's own dataset is not per se a conflict; any required COI disclosure is handled at submission.
%% Trimmed + bold run-in labels removed 2026-06-16 (rebalance vs the expanded framework section).
%% Summarised 2026-06-17 to 3 paragraphs; PT-CS para now LEADS WITH CONTAINMENT (framework rests on the
%% public datasets; PT-CS only for transfer + the PT-CS-only experiments; verified stratum; conversation
%% gold-independent) so the gold limits do not read as undermining the article.
%% Inflation-proxy digits regenerate via phase5.signed_deviation on the declared baseline cell.

\section{Threats to Validity}
\label{sec:threats}

The PT-CS dataset was collected by the authors of this article and carries real limits: its grades come
from a mix of sources of uneven reliability (a requalification course below university level, not every
grade teacher-reviewed), the curation criterion shows intervention rather than correctness, no per-teacher scores
allow an inter-rater check, and an AI-seeded suggestion could inflate agreement if its model resembled one
we test. These limitations keep this
dataset from reaching the paper's conclusions. The framework rests on the three public datasets, whose
references are independent of PT-CS. PT-CS enters only where nothing else can serve, as the transfer test
and in the experiments that need its multi-question submissions: grading scope, criterion decomposition,
and the conversation sub-study. There, every agreement result uses the teacher-reviewed
verified stratum (Section~\ref{sec:datasets}), and the conversation sub-study measures model behaviour
without depending on the gold at all. A lenient course
standard could also inflate the grades, but the model-minus-gold deviation does not clearly bear this out.

The results are conditional on the served backends, so the transferable claim is the configuration
guidance, not the absolute scores. Temperature 0 is not deterministic: 10 to 43\% of items still vary
across repetitions (Section~\ref{sec:discussion}). Reasoning cost is comparable only within a backend, not
across them. The reasoning benefit is established on the tractable items only, since the non-random
truncation exclusion drops the longest-reasoning cases.

The guidance was established on the model families tested, with reasoning isolated within family, and is
keyed on observable properties rather than model identity, so its reach to other models is by design rather
than proof. Across all experiments, code evidence is narrower than short-answer evidence, resting on two
Java-centric datasets, so we draw the code conclusions more cautiously. GPT-5.1 is a reference point, not an
inference target: it graded small subsets and corroborates rather than ranks. The guidance contrast bundles
two different interventions, a rubric for code and a reference answer for short answers, which we report
separately rather than as one combined effect. The expensive factors were varied one at a time rather than
fully crossed, so interactions between them, such as reasoning with decomposition, are unmeasured.
